\documentclass[10pt,a4paper]{amsart}
\usepackage[margin=19mm]{geometry}
\usepackage{amsmath,amssymb,mathtools}
\usepackage[scr=rsfs,cal=euler]{mathalfa}
\usepackage{xfrac}
\usepackage[unicode,hidelinks,bookmarks=false]{hyperref}
\newcommand{\ie}{i.e.\ }
\newcommand{\cf}{cf.\ }
\newcommand{\resp}{respectively\ }
\newcommand{\Subsection}{\subsection}
\providecommand{\nobreakdash}{\nobreak\hbox{-}}
\setlength{\emergencystretch}{2em}
\multlinegap0pt
\usepackage{amssymb}
\usepackage{esint}
\newtheorem{prop}{Proposition}
\newtheorem{lem}{Lemma}
\title{Constructing $\lambda$-Angenent curve by flow method}
\author{Pak Tung Ho}
\date{Source: arXiv:2601.02853v1 (CC BY 4.0)}
\begin{document}
\maketitle

\noindent\emph{Source and licence.} Constructing $\lambda$-Angenent curve by flow method. Version arXiv:2601.02853v1 (CC BY 4.0), distributed by arXiv under the Creative Commons Attribution 4.0 International licence, http://creativecommons.org/licenses/by/4.0/, as recorded in the arXiv licence metadata for this exact version and shown on the abstract page https://arxiv.org/abs/2601.02853v1. Full licence notice: \texttt{licenses/arxiv-2601.02853-CC-BY-4.0.txt}.\par\smallskip

\noindent\textit{Editorial scope.} Counted unit: the article's prop prop9 (source0.tex lines 448--452). The article's complete original proof of it is delivered with it (source0.tex lines 453--510, one \texttt{proof} environment of 58 lines). No mathematics is written, repaired or re-derived: the statement and the proof are the article's own lines, in the article's own notation, and no retained line is substituted or re-flowed.  Retained uncounted as the source-local context the proof needs: lines 233--239; lines 431--435; lines 439--441.  Nothing was omitted from any retained span, so the package carries no omission record.  No retained line contains a citation, so the wrapper carries no bibliography and no bibliography entry.  No retained line contains a figure environment, an image-inclusion command or drawing code, so no image is delivered and the package declares no asset dependency.  Editorial formatting only, in addition: an AMS \texttt{amsart} preamble that restates the article's own declarations and macros used by the retained lines, namely the environments \texttt{prop}, \texttt{lem} on separate counters, together with the standard packages those lines need.  The retained environments print in this document as Proposition 1, Lemma 1 in order of appearance, whereas the article prints the same items with the numbering of its own full document.  The complete native source of the article as distributed in the arXiv e-print for this exact version is bundled unchanged as \texttt{sources/192187/source0.tex}.
\begin{equation}\label{1.5}
\begin{split}
L_g(\mathcal{C})&=\int_{-\infty}^\infty (2(\lambda-1))^{\frac{\lambda-1}{2}}e^{-\frac{2(\lambda-1)+u^2}{4}}du
=2\sqrt{\pi}\left(\frac{2(\lambda-1)}{e}\right)^{\frac{\lambda-1}{2}},\\
L_g(\mathcal{P})&=\int_0^\infty u^{\lambda-1} e^{-\frac{u^2}{4}}du.
\end{split}
\end{equation}

\begin{equation}\label{2.1}
L_g[a,b,c]
:=2(a^{\lambda-1} e^{-a^2/4}+b^{\lambda-1} e^{-b^2/4})\int_0^ce^{-x^2/4}dx
+2e^{-c^2/4}\int_a^b r^{\lambda-1} e^{-r^2/4}dr.
\end{equation}

\begin{equation}\label{2.2}
\frac{e^{-c_0^2/4}}{\int_0^{c_0}e^{-x^2/4}dx}=2.
\end{equation}

\begin{prop}\label{prop9}
Suppose that $\lambda>1$.
For every $a,b\in (0,\infty)$, we have
$$L_g[a,b,c]<2L_g(\mathcal{P}).$$
\end{prop}

\begin{proof}
In view of (\ref{2.2}), we can write (\ref{2.1}) as
$$L_g(a,b,c_0)=f_\lambda(a)+g_\lambda(b),$$
where $M=\displaystyle\int_0^{c_0}e^{-x^2/4}dx$ and
\begin{equation*}
\begin{split}
f_\lambda(a)&=2a^{\lambda-1} e^{-a^2/4}M+4M\int_a^0r^{\lambda-1} e^{-r^2/4}dr,\\
g_\lambda(b)&=2b^{\lambda-1} e^{-b^2/4}M+4M\int_0^br^{\lambda-1} e^{-r^2/4}dr.
\end{split}
\end{equation*}
The derivative of $f$ is
$$f_\lambda'(a)=-a^{\lambda-2}e^{-a^2/4}M\big(a^2+4a-2(\lambda-1)\big).$$
Therefore, $\displaystyle\max_{a\in[0,\infty)}f_\lambda(a)$
is achieved at $a_\lambda=-2+\sqrt{2(1+\lambda)}$.
Similarly, one can show that
$\displaystyle\max_{b\in[0,\infty)}g_\lambda(b)$
is achieved at $b_\lambda=2+\sqrt{2(1+\lambda)}$.
It follows from (\ref{1.5}) that
$$L_g(\mathcal{P})=\int_0^\infty u^{\lambda-1} e^{-\frac{u^2}{4}}du.$$
Hence, it suffices to show that
\begin{equation}\label{2.4}
f_\lambda(a_{\lambda})+g_\lambda(b_\lambda)<2\int_0^\infty r^{\lambda-1} e^{-\frac{r^2}{4}}dr
\end{equation}
for any $\lambda>1$.

To prove (\ref{2.4}), we are going to prove the following two lemmas:


\begin{lem}\label{induction_lem}
The inequality \eqref{2.4} holds whenever $1\leq \lambda\leq 7$.
\end{lem}


\begin{lem}\label{lem10} There holds
\begin{equation}\label{2.3}
f_{\lambda+2}(a_{\lambda+2})+g_{\lambda+2}(b_{\lambda+2})
<2\lambda\big(f_{\lambda}(a_{\lambda})+g_{\lambda}(b_{\lambda})\big)
\end{equation}
for all $\lambda>\lambda_0$.
\end{lem}

Let us assume Lemma \ref{induction_lem} and Lemma \ref{lem10} for the time being;
and we can see that (\ref{2.4}) holds for all $\lambda>1$.
It follows from integration by parts that
\begin{equation*}
\begin{split}
\int_0^\infty r^{\lambda+1} e^{-\frac{r^2}{4}}dr
&=\Big[-2r^{\lambda}e^{-\frac{r^2}{4}}\Big]_0^\infty+2\lambda\int_0^\infty r^{\lambda-1} e^{-\frac{r^2}{4}}dr\\
&=2\lambda\int_0^\infty r^{\lambda-1} e^{-\frac{r^2}{4}}dr.
\end{split}
\end{equation*}
This together with Lemma \ref{lem10} implies that
if (\ref{2.4}) holds for a particular value of $\lambda>\lambda_0$, then
(\ref{2.4}) holds for $\lambda+2$.
Since Lemma \ref{induction_lem} asserts that (\ref{2.4}) holds for all $1\leq \lambda\leq 7$,
we can conclude that  (\ref{2.4}) holds for all $\lambda>1$.
This proves Proposition \ref{prop9} by assuming Lemma \ref{induction_lem} and Lemma \ref{lem10}.
\end{proof}

\end{document}
